\section{收获与心得}
通过本次实验,对操作系统中文件系统的实现原理和并发控制机制有了更深入的理解,主要收获如下:

\textbf{1. 文件系统底层实现的理解。} 在任务一中,通过将内存数组模拟为硬盘,亲身实践了文件系统从底层数据结构到上层接口的完整构建过程。FAT表的链式管理方式直观展示了磁盘空间分配与回收的机制,目录项结构体揭示了文件元数据管理的基本方法,文件打开表的引入则体现了操作系统通过缓存机制提升文件访问效率的设计思路。这些实践使课本中抽象的理论概念变得具体可感。

\textbf{2. 读者-写者问题的工程实现。} 任务二将经典的读者-写者同步问题应用于实际文件系统场景中。实现过程中需仔细考虑互斥锁和信号量的配合:用mutex保护reader\_count的修改、用rw\_lock信号量控制资源访问权限。这加深了对P/V操作语义的理解,尤其是第一个读者和最后一个读者在资源信号量上的"接力"逻辑。此外,使用\texttt{sem\_trywait}替代\texttt{sem\_wait}以解决交互界面响应问题,也体现了理论算法需结合实际场景灵活调整的工程思维。
% \end{content}
